% GENERATED FILE. Edit canonical lessons, then run npm run generate:books.
% canonical-source-hash: fnv1a64:d8c619018a1490e4
% canonical-lessons: ZH-C272-qizi, ZH-C272-laogong, ZH-W272-lao-old, ZH-C272-waigong, ZH-C272-shushu, ZH-C272-ayi

\chapter{Husband and Wife}
\label{ch:zh-husband-and-wife}
\hlchaptermodality{272}

\begin{chapteropening}
\textbf{By the end of this chapter:} \emph{I can talk about a husband and a wife, my grandfather on my mother's side, and an uncle and an auntie.}

Complete the last lesson of chapter 272: I can talk about a husband and a wife, my grandfather on my mother's side, and an uncle and an auntie.
\end{chapteropening}

\section[qīzi]{qīzi — wife}

\label{lesson:ZH-C272-qizi}



\begin{quote}
Before the new one: say the Mandarin for bored, then the Mandarin for something on your mind.
\end{quote}

\subsection*{You'll want to know: qīzi}
\emph{qīzi} — \textquotedblleft{}wife\textquotedblright{}.

\begin{quote}
\emph{tā de qīzi shì yīshēng} — His wife is a doctor.
\end{quote}

Its characters come later; for now it is a word for the ear.

\begin{cousinweb}
\emph{qīzi} is built from \emph{qī} and \zh{子}, both inherited from Middle Chinese, \emph{tshej} and \emph{tsiX} (Guangyun, in Baxter's transcription), which glosses the first 'equal'. In the word the second loses its tone.
\end{cousinweb}

\begin{grammarlens}[title={What you've built}]
One more word for what I have and can do.
\end{grammarlens}

\subsection*{Your turn}
\begin{itemize}
\raggedright
  \item \emph{Say it:} \emph{qīzi}
  \item \emph{Say it:} \emph{qīzi}, once more
  \item \emph{Say it:} \emph{tā de qīzi shì yīshēng} — His wife is a doctor.
  \item \emph{Recall:} say the Mandarin for bored, then the Mandarin for something on your mind, then say \emph{qīzi} again
\end{itemize}

\begin{tcolorbox}[breakable,colback=teal!4,colframe=teal!35!black,title={Before you move on}]
What does qīzi mean? (\textquotedblleft{}Wife\textquotedblright{}.) Say it once more.
\end{tcolorbox}
\section[\texorpdfstring{\zh{老公} (lǎogōng)}{老公 (lǎogōng)}]{\zh{老公} (lǎogōng) — hubby; husband (informal)}

\label{lesson:ZH-C272-laogong}



\begin{quote}
Before the new one: say the Mandarin for something on your mind, then the Mandarin for a wife.
\end{quote}

\subsection*{You'll want to know: \zh{老公}}
\textbf{\zh{老公}} — \emph{lǎogōng} — \textquotedblleft{}hubby; husband (informal)\textquotedblright{}.

\begin{quote}
\textbf{\zh{我老公先回家了}} \emph{wǒ lǎogōng xiān huí jiā le} — My husband has gone home ahead of me.
\end{quote}

The new character, in pieces: \zh{老} — six strokes: a short bar, a vertical and a long bar, then a long stroke falling to the left; under them a short stroke falling to the left and a vertical that bends right and hooks up, the right half of \zh{比}. The next lesson puts it in your hand.

\begin{cousinweb}
\textbf{\zh{老公}} (\emph{lǎogōng}) is built from \zh{老} 'old' and \zh{公}, both inherited from Middle Chinese, \emph{lawX} and \emph{kuwng} (Guangyun, in Baxter's transcription), which glosses the first 'aged' and the second, among other senses, 'a father' and 'an official'. How the pair came to mean a husband is not settled.
\end{cousinweb}

\begin{grammarlens}[title={What you've built}]
One more word for what I have and can do.
\end{grammarlens}

\subsection*{Your turn}
\begin{itemize}
\raggedright
  \item \emph{Say it:} \emph{lǎogōng}
  \item \emph{Say it:} \emph{lǎogōng}, once more
  \item \emph{Say it:} \emph{wǒ lǎogōng xiān huí jiā le} — My husband has gone home ahead of me.
  \item \emph{Recall:} say the Mandarin for something on your mind, then the Mandarin for a wife, then say \emph{lǎogōng} again
\end{itemize}

\begin{tcolorbox}[breakable,colback=teal!4,colframe=teal!35!black,title={Before you move on}]
What does \zh{老公} mean? (\textquotedblleft{}Hubby; husband (informal)\textquotedblright{}.) Say it once more.
\end{tcolorbox}
\section[\texorpdfstring{\zh{老} (lǎo)}{老 (lǎo)}]{\zh{老} — a cross and a long fall, over the right half of \zh{比}}

\label{lesson:ZH-W272-lao-old}



\begin{quote}
The word first, then the character inside it.

\begin{itemize}
\raggedright
  \item \emph{Recall:} say \emph{lǎogōng} — \textbf{R1}, one lesson back, before it decays
\end{itemize}

\begin{quote}
\zh{比} \zh{公}
\end{quote}

\emph{Read it:} each one aloud, with nothing said first — \emph{bǐ}, \emph{gōng}
\end{quote}

\subsection*{Script — a cross and a long fall, over the right half of \zh{比}}
\hlblockfigure{\detokenize{figures/ZH-W272-lao-old-filmstrip.pdf}}{How \zh{老} is written, stroke by stroke}

\begin{quote}
\zh{老}
\end{quote}

\textbf{Two pieces, top before bottom.}

\textbf{The top piece, four strokes.} A short bar. A vertical down through it. A long bar under them. Then a long stroke from the upper right falling all the way to the lower left.

\textbf{The bottom piece, two strokes,} the right half of \zh{比}: a short stroke falling to the left, then a vertical that bends to the right along the foot and hooks up.

4 + 2 = six.

\subsection*{Your turn}
\begin{itemize}
\raggedright
  \item \emph{Write it:} \zh{老}, all six, with the model in front of you
  \item \emph{Cover:} write it again from memory, and count the pieces, not the strokes
\end{itemize}

\begin{tcolorbox}[breakable,colback=teal!4,colframe=teal!35!black,title={Before you move on}]
How many strokes in \zh{老}? (\textbf{Six}.) Which stroke runs across the whole character? (\textbf{The long fall to the left}.) Where have you made the last two strokes? (\textbf{On the right of \zh{比}}.)

Source: \href{https://raw.githubusercontent.com/chanind/hanzi-writer-data/68d10a4b21150cae5e1ebbd223eed289cf32d90c/data/\%E8\%80\%81.json}{Hanzi Writer Data \zh{老}, pinned stroke paths and medians}.
\end{tcolorbox}
\section[\texorpdfstring{\zh{外公} (wàigōng)}{外公 (wàigōng)}]{\zh{外公} (wàigōng) — grandpa on the mother's side}

\label{lesson:ZH-C272-waigong}



\begin{quote}
Before the new one: say the Mandarin for a wife, then the Mandarin for hubby.
\end{quote}

\subsection*{You'll want to know: \zh{外公}}
\textbf{\zh{外公}} — \emph{wàigōng} — \textquotedblleft{}grandpa on the mother's side\textquotedblright{}.

\begin{quote}
\textbf{\zh{外公很热心}} \emph{wàigōng hěn rèxīn} — My grandpa is very warm-hearted.
\end{quote}

\begin{cousinweb}
\textbf{\zh{外公}} (\emph{wàigōng}) is built from \zh{外} 'outside' and \zh{公} 'a father', both inherited from Middle Chinese, \emph{ngwajH} and \emph{kuwng} (Guangyun, in Baxter's transcription), which glosses the second, among other senses, 'a father'.
\end{cousinweb}

\begin{grammarlens}[title={What you've built}]
One more word for what I have and can do.
\end{grammarlens}

\subsection*{Your turn}
\begin{itemize}
\raggedright
  \item \emph{Say it:} \emph{wàigōng}
  \item \emph{Say it:} \emph{wàigōng}, once more
  \item \emph{Say it:} \emph{wàigōng hěn rèxīn} — My grandpa is very warm-hearted.
  \item \emph{Recall:} say the Mandarin for a wife, then the Mandarin for hubby, then say \emph{wàigōng} again
\end{itemize}

\begin{tcolorbox}[breakable,colback=teal!4,colframe=teal!35!black,title={Before you move on}]
What does \zh{外公} mean? (\textquotedblleft{}Grandpa on the mother's side\textquotedblright{}.) Say it once more.
\end{tcolorbox}
\section[shūshu]{shūshu — uncle; father's younger brother}

\label{lesson:ZH-C272-shushu}



\begin{quote}
Before the new one: say the Mandarin for hubby, then the Mandarin for grandpa on the mother's side.
\end{quote}

\subsection*{You'll want to know: shūshu}
\emph{shūshu} — \textquotedblleft{}uncle; father's younger brother\textquotedblright{}.

\begin{quote}
\emph{wǒ shūshu shì diàngōng} — My uncle is an electrician.
\end{quote}

Its characters come later; for now it is a word for the ear.

\begin{cousinweb}
\emph{shūshu} is built from \emph{shū}, said twice, the second time with no tone. It is inherited from Middle Chinese \emph{syuwk} (Guangyun, in Baxter's transcription), glossed 'a father's youngest brother'.
\end{cousinweb}

\begin{grammarlens}[title={What you've built}]
One more word for what I have and can do.
\end{grammarlens}

\subsection*{Your turn}
\begin{itemize}
\raggedright
  \item \emph{Say it:} \emph{shūshu}
  \item \emph{Say it:} \emph{shūshu}, once more
  \item \emph{Say it:} \emph{wǒ shūshu shì diàngōng} — My uncle is an electrician.
  \item \emph{Recall:} say the Mandarin for hubby, then the Mandarin for grandpa on the mother's side, then say \emph{shūshu} again
\end{itemize}

\begin{tcolorbox}[breakable,colback=teal!4,colframe=teal!35!black,title={Before you move on}]
What does shūshu mean? (\textquotedblleft{}Uncle; father's younger brother\textquotedblright{}.) Say it once more.
\end{tcolorbox}
\section[āyí]{āyí — auntie; mother's sister}

\label{lesson:ZH-C272-ayi}



\begin{quote}
Before the new one: say the Mandarin for grandpa on the mother's side, then the Mandarin for an uncle.
\end{quote}

\subsection*{You'll want to know: āyí}
\emph{āyí} — \textquotedblleft{}auntie; mother's sister\textquotedblright{}.

\begin{quote}
\emph{nàge āyí shì wǒ māma de péngyou} — That lady is a friend of my mum's.
\end{quote}

Its characters come later; for now it is a word for the ear.

\begin{cousinweb}
\emph{āyí} is built from \emph{ā}, a friendly prefix, and \emph{yí} 'a mother's sister', both inherited from Middle Chinese, \emph{'a} and \emph{yij} (Guangyun, in Baxter's transcription), which glosses the second 'a mother's sister'. The dictionary gives the first only as 'a bend' and 'a hill', so its use before names is not settled.
\end{cousinweb}

\begin{grammarlens}[title={What you've built}]
One more word for what I have and can do.
\end{grammarlens}

\subsection*{Your turn}
\begin{itemize}
\raggedright
  \item \emph{Say it:} \emph{āyí}
  \item \emph{Say it:} \emph{āyí}, once more
  \item \emph{Say it:} \emph{nàge āyí shì wǒ māma de péngyou} — That lady is a friend of my mum's.
  \item \emph{Recall:} say the Mandarin for grandpa on the mother's side, then the Mandarin for an uncle, then say \emph{āyí} again
\end{itemize}

\begin{tcolorbox}[breakable,colback=teal!4,colframe=teal!35!black,title={Before you move on}]
What does āyí mean? (\textquotedblleft{}Auntie; mother's sister\textquotedblright{}.) Say it once more.
\end{tcolorbox}
